% TOP_PROBLEM: {"id":30007663,"problem_number":"LOCAL-30007663","title":"Building accountable state capacity","category_id":21,"category":{"id":21,"name":"economics","display_name":"Economics","description":"Open economic theory statements, published research agendas, and proposed scoped specifications; question provenance and historical priority are qualified per record.","slug":"economics","order_index":21,"created_at":"2026-10-08T00:00:00Z"},"status":"research_question","external_url":null,"legacy_ids":[],"published":false,"statement":"In the explicitly specified setting: Which sequences of tax, legal, and administrative reforms build a capable state while preserving citizen accountability and limiting elite appropriation or repression? The mathematical statement above fixes the target, assumptions and scope of this catalogue formulation.","economics":{"original_id":152,"field":"Political economy and institutions","kind":"design","formalization_basis":"adapted_research_question","source_status":"research_agenda","priority_status":"earliest_found","priority_note":"Earliest directly inspected qualifying passage in this audit; earlier origins were not established. A review or research-agenda statement does not imply original priority.","earliest_verified_reference":{"authors":["Timothy Besley","Robin Burgess","Adnan Khan","Guo Xu"],"title":"Bureaucracy and Development","year":2022,"url":"https://www.robinburgess.com/s/Bureaucracy_and_Development.pdf","doi":"10.1146/annurev-economics-080521-011950","locator":"Section 4.2, printed pp. 413–415; concluding research agenda, printed p. 418, third direction","evidence_summary":"The review calls for separating politician and bureaucrat roles and learning how political accountability shapes bureaucratic effectiveness."},"current_reference":{"authors":["Timothy Besley","Robin Burgess","Adnan Khan","Guo Xu"],"title":"Bureaucracy and Development","year":2022,"url":"https://www.robinburgess.com/s/Bureaucracy_and_Development.pdf","doi":"10.1146/annurev-economics-080521-011950","locator":"Section 4.2, printed pp. 413–415; concluding research agenda, printed p. 418, third direction","evidence_summary":"The review calls for separating politician and bureaucrat roles and learning how political accountability shapes bureaucratic effectiveness."},"formalization_note":"The cited passage establishes a research direction, not this exact mathematical statement. The notation, population, policy menu, assumptions, and estimands are an original adaptation for this catalogue; no universal unsolved theorem or source priority is claimed."},"statusQualification":"Published question or research agenda verified at the cited locator. The mathematical specification is adapted for this catalogue and includes additional assumptions; source publication does not establish first-ever priority or certify this exact model remains unsolved. The cited passage establishes a research direction, not this exact mathematical statement. The notation, population, policy menu, assumptions, and estimands are an original adaptation for this catalogue; no universal unsolved theorem or source priority is claimed.","statusReviewedAt":"2026-10-08"}
% FIRST_POSTED: 2026-10-08
% ENTRY_KIND: statement_only
% !TeX program = lualatex
\documentclass[11pt]{article}
\usepackage{fontspec}
\usepackage[a4paper,margin=25mm]{geometry}
\usepackage{amsmath,amssymb,mathtools}
\usepackage{xurl}
\usepackage[hidelinks]{hyperref}
\newcommand{\sref}[2]{\hyperlink{source:#1:#2}{[S#2]}}
\setlength{\emergencystretch}{3em}
\begin{document}
% problemId:30007663
\section*{Building accountable state capacity}\label{30007663}
\subsection{Definitions and mathematical statement}
Fix polities initially possessing limited fiscal, legal, and administrative capacity. Let \(p\) be a feasible sequence of capacity and accountability reforms with a stated cost and implementation horizon. Let \(K_{ct}(p)\) be verified public-service and revenue capacity, \(A_{ct}(p)\) an accountability index defined before evaluation, and \(R_{ct}(p)\) measured coercion or elite appropriation. For horizon \(T\), estimate \(\tau_K(p)=E[K_{cT}(p)-K_{cT}(0)]\), and analogous accountability and coercion effects. The design objective is to find sequences improving capacity subject to specified lower bounds on accountability and upper bounds on repression and appropriation. These constraints encode declared social objectives. Separate the roles of elected politicians, appointed officials, and nonstate actors, because identical administrative investments may operate under different political incentives. Identification requires credible reform timing or assignment and explicit treatment of conflict, spillovers, and regime change. Evaluate service outcomes and citizens' ability to contest decisions rather than equating greater extraction with better capacity. The review explicitly requests knowledge about political oversight and bureaucratic effectiveness. Sequencing, causal estimands, and the constrained reform objective are our adaptation; the source does not offer a universally valid reform recipe.

\par\textbf{Formulation and scope.} The cited passage establishes a research direction, not this exact mathematical statement. The notation, population, policy menu, assumptions, and estimands are an original adaptation for this catalogue; no universal unsolved theorem or source priority is claimed.

\par\textbf{Open-question provenance.} An explicit question or research agenda was verified at the source locator. This does not by itself certify that every later version remains unresolved \sref{30007663}{1}.

\par\textbf{Historical priority.} Earliest directly inspected qualifying passage in this audit; earlier origins were not established. A review or research-agenda statement does not imply original priority.

\subsection{Short English statement}
In the explicitly specified setting: Which sequences of tax, legal, and administrative reforms build a capable state while preserving citizen accountability and limiting elite appropriation or repression? The mathematical statement above fixes the target, assumptions and scope of this catalogue formulation.

\subsection{Sources}
\begin{itemize}\raggedright
\item[\textnormal{[S1]}]\hypertarget{source:30007663:1}{} Timothy Besley, Robin Burgess, Adnan Khan, Guo Xu. 2022. Bureaucracy and Development. Section 4.2, printed pp. 413–415; concluding research agenda, printed p. 418, third direction.
\url{https://www.robinburgess.com/s/Bureaucracy_and_Development.pdf}.
DOI: 10.1146/annurev-economics-080521-011950.
{\small\textit{Earliest verified question/agenda reference; historical priority qualified below. The review calls for separating politician and bureaucrat roles and learning how political accountability shapes bureaucratic effectiveness.}}
\end{itemize}
\end{document}
